\thispagestyle{plain}
\begin{center}
	\vspace*{1.5cm}
	{\Large \bfseries  Abstract}
\end{center}
\vspace{0.5cm}
This research aimed to implement time series and Machine Learning models to improve decision-making in the supplier replenishment process of Donna S.A.C. The problem addressed is related to the difficulty of adequately estimating product demand, especially in contexts where sales variability, seasonality, and the short shelf life of certain products may lead to overstock, stockouts, economic losses, and purchasing decisions based mainly on the administrator's experience.

The research was developed under an applied and quantitative approach, using the company's historical sales data as the main source of information. The methodological process included data acquisition and understanding, information preprocessing, selection of relevant variables, construction of temporal features, training of predictive models, comparative evaluation, and deployment of the solution. During preprocessing, activities such as data cleaning, outlier treatment, sales aggregation by date, selection of relevant columns, and generation of variables such as day of the week, month, sales lags, and moving averages were carried out.

To evaluate model performance, metrics such as MAE, MSE, RMSE, MAPE, and the coefficient of determination $R^2$ were used, with the purpose of comparing the generated predictions against the actual values. The results obtained made it possible to identify the models that presented the best overall performance for each category, as they recorded lower error levels and a better fit compared to the other evaluated models. Likewise, a support solution was developed through a system that allows users to visualize predictions, compare actual and forecasted sales, and facilitate replenishment planning.

Finally, it is concluded that the implementation of predictive models contributes to transforming historical sales data into useful information for operational decision-making. The proposed solution helps reduce uncertainty in product purchasing, optimize the supply budget, and reduce risks associated with excess or insufficient inventory, thereby strengthening the efficiency of the supplier ordering process at Donna S.A.C.


\textbf{Keywords: } Forecasting, Sales, Machine Learning, Regression